% !TEX program = pdflatex
\documentclass[11pt, a4paper]{article}

% ── Packages ──────────────────────────────────────────────────────────────────
\usepackage[utf8]{inputenc}
\usepackage[T1]{fontenc}
\usepackage{lmodern}
\usepackage[margin=1in]{geometry}
\usepackage{graphicx}
\usepackage{booktabs}
\usepackage{hyperref}
\usepackage{xcolor}
\usepackage{amsmath}
\usepackage{enumitem}
\usepackage{float}
\usepackage{tikz}
\usepackage{listings}
\usepackage[numbers]{natbib}
\usepackage{microtype}
\usepackage{abstract}

% ── Macros ────────────────────────────────────────────────────────────────────
\newcommand{\sprintagent}{\textsc{SprintAgent}}
\newcommand{\ioc}{Inversion of Control}
\newcommand{\di}{Dependency Injection}

% ── Listings config ──────────────────────────────────────────────────────────
\lstset{
  basicstyle=\ttfamily\small,
  keywordstyle=\bfseries\color{blue!70!black},
  commentstyle=\color{green!50!black},
  stringstyle=\color{red!60!black},
  breaklines=true,
  frame=single,
  columns=fullflexible,
  keepspaces=true,
  showstringspaces=false,
  captionpos=b,
}

% ── Hyperref config ──────────────────────────────────────────────────────────
\hypersetup{
  colorlinks=true,
  linkcolor=blue!70!black,
  citecolor=blue!70!black,
  urlcolor=blue!70!black,
}

% ══════════════════════════════════════════════════════════════════════════════
\title{%
  \sprintagent{}: A Production-Ready Framework for\\
  Hierarchical Autonomous Agentic Systems%
}
\author{Ali Shamsaddinlou}
\date{March 2026}

\begin{document}
\maketitle

% ── Abstract ─────────────────────────────────────────────────────────────────
\begin{abstract}
From an engineering perspective, building autonomous agentic systems
with large language models (LLMs) demands more deliberate design
thinking than traditional software.  The trade-offs are sharper and
the consequences of getting them wrong are harder to debug: complexity
vs.\ explainability, isolation vs.\ throughput, context growth vs.\
performance, and configuration flexibility vs.\ simplicity.  Yet most
prototypes lack the engineering infrastructure needed to navigate these
trade-offs in production.  This paper introduces \sprintagent{}, an
open-source framework that provides a production-ready scaffold for
hierarchical autonomous agentic systems, developed with these
engineering trade-offs explicitly in mind so that the resulting systems
remain manageable, auditable, and explainable.  The framework features
a planner/designer/critic orchestration loop per agent, Hydra-based
\ioc{} configuration, SQLite-backed session management with turn
trimming, checkpoint-based auto-reset on performance degradation,
isolated parallel execution, and unified logging for observability.
The framework is designed to be highly modular: practitioners can plug
their own logic into each module and extend on top with minimal
friction.  To demonstrate how this looks in practice, an agentic
project-management system is built on top of \sprintagent{} as a
reference application; the project-management system is not the
contribution itself but rather a demonstration of the scaffold in
action.
\end{abstract}

% ── 1. Introduction ──────────────────────────────────────────────────────────
\section{Introduction}
\label{sec:introduction}

Large language models have enabled a new class of autonomous software
agents that can plan, reason, and execute multi-step tasks with minimal
human intervention~\cite{wang2024survey, xi2023rise}.  Frameworks such as
LangChain~\cite{chase2023langchain}, AutoGen~\cite{wu2023autogen},
MetaGPT~\cite{hong2023metagpt}, and CrewAI~\cite{crewai2024} have
lowered the barrier to building agentic applications.  However, moving
from a notebook prototype to a production-grade agentic system reveals
a gap: most existing frameworks focus on the \emph{agent logic} (prompt
chaining, tool calling, memory) while leaving the surrounding
\emph{engineering infrastructure}---configuration management, fault
tolerance, session persistence, observability, and execution
isolation---to the developer.

From an engineering perspective, building agentic systems demands more
deliberate design thinking than traditional software.  The trade-offs
are sharper and the consequences of getting them wrong are harder to
debug:
\begin{itemize}[leftmargin=*]
  \item \textbf{Complexity vs.\ explainability} -- powerful multi-agent
    orchestration is useless if no one can understand what the system
    did and why.
  \item \textbf{Isolation vs.\ throughput} -- strong fault boundaries
    keep runs safe but add overhead.
  \item \textbf{Context growth vs.\ performance} -- agents accumulate
    history that must be managed, not just appended.
  \item \textbf{Configuration flexibility vs.\ simplicity} -- real-world
    systems have dozens of knobs; the config system must scale without
    becoming its own problem.
\end{itemize}

What does a production-ready autonomous agentic system look like when it
needs to handle complex agents without becoming a mess?  This question
motivated the design of \sprintagent{}, an open-source framework
developed with these engineering trade-offs explicitly in mind.
\sprintagent{} is a highly modular scaffold where practitioners can plug
their own logic into each module and extend on top.  For demonstration
purposes, an agentic project-management system is built on top of the
framework to show what it looks like in practice to build such a system.

The key contributions of this paper are:
\begin{enumerate}[leftmargin=*]
  \item A \textbf{hierarchical agentic architecture} in which each stage
    (project, module, task, visualization) is handled by an autonomous
    agent that internally runs a planner/designer/critic refinement loop.
  \item An \textbf{\ioc{} configuration system} built on Hydra and
    OmegaConf with custom resolvers, enabling new agents and experiments
    to be wired entirely through YAML without code changes.
  \item \textbf{Production infrastructure} including SQLite-backed
    session management with turn trimming and summarization, checkpoint-based
    auto-reset on quality degradation, process-isolated parallel
    execution, and unified file-and-terminal logging for observability
    and explainability.
  \item An \textbf{open-source implementation} accompanied by a
    demonstration application---an agentic project-management
    system---that illustrates how to build a complete hierarchical
    agentic application on the framework.
\end{enumerate}

% ── 2. Background ────────────────────────────────────────────────────────────
\section{Background and Motivation}
\label{sec:background}

\subsection{LLM-Based Autonomous Agents}

The ReAct paradigm~\cite{yao2023react} showed that interleaving
reasoning and action in LLMs produces more capable agents.
Reflexion~\cite{shinn2023reflexion} added verbal self-reflection,
enabling agents to learn from mistakes within a single episode.  These
ideas have been operationalized in multi-agent frameworks: AutoGen
introduces conversable agents that collaborate through
messages~\cite{wu2023autogen}; MetaGPT assigns software-engineering
roles to agents in a simulated company~\cite{hong2023metagpt}; CrewAI
orchestrates role-playing agents with configurable
workflows~\cite{crewai2024}; and the OpenAI Agents SDK provides
primitives for tool-calling agents with handoff
support~\cite{openaiagentssdk2025}.

\subsection{The Production Gap}

While these frameworks excel at composing agent logic, production
deployments require additional concerns:
\begin{itemize}[leftmargin=*]
  \item \textbf{Configuration complexity.}  Real-world agentic systems
    have dozens of knobs (model selection, temperature, scoring
    thresholds, memory policies, pipeline stages).  Hard-coded or
    flat-file configs do not scale.
  \item \textbf{Fault tolerance.}  LLM outputs are non-deterministic;
    iterative refinement can degrade quality.  Systems need mechanisms
    to detect and recover from regressions.
  \item \textbf{Session persistence and memory management.}  Long-running
    agents accumulate context that exceeds token limits.  Persistent
    sessions with turn trimming and summarization are essential.
  \item \textbf{Observability.}  Debugging multi-agent systems requires
    unified logging that captures what each agent did, when, and why.
  \item \textbf{Execution isolation.}  Parallel runs must not interfere
    with each other; failures in one run should not cascade.
\end{itemize}

\sprintagent{} addresses each of these concerns as a first-class
framework feature, so developers can focus on domain logic rather than
infrastructure.

% ── 3. System Overview ───────────────────────────────────────────────────────
\section{System Overview}
\label{sec:overview}

\sprintagent{} transforms a high-level project brief into an
execution-ready roadmap through a four-stage hierarchical pipeline:
\textbf{project} $\rightarrow$ \textbf{module} $\rightarrow$
\textbf{task} $\rightarrow$ \textbf{visualization}.  Each stage is
handled by a specialized agent that refines its output through an
internal planner/designer/critic loop before passing results to the
next level.

\begin{figure}[H]
  \centering
  \includegraphics[width=\textwidth]{figures/pipeline.png}
  \caption{Hierarchical pipeline.  Each stage consumes the output of the
    previous stage and produces a refined plan.  The visualization stage
    generates Mermaid diagrams rendered to PNG/PDF.}
  \label{fig:pipeline}
\end{figure}

Table~\ref{tab:stages} summarizes the inputs, outputs, and responsible
agent for each stage.

\begin{table}[H]
  \centering
  \caption{Pipeline stages and their agents.}
  \label{tab:stages}
  \begin{tabular}{@{}llll@{}}
    \toprule
    \textbf{Stage} & \textbf{Input} & \textbf{Output} & \textbf{Agent} \\
    \midrule
    Project       & User brief           & \texttt{project\_plan.md}  & StatefulProjectAgent \\
    Module        & Project plan         & \texttt{module\_plan.md}   & StatefulModuleAgent  \\
    Task          & Module plan          & \texttt{task\_plan.md}     & StatefulTaskAgent    \\
    Visualization & All three plans      & Mermaid + PNG/PDF diagrams & StatefulVisualizationAgent \\
    \bottomrule
  \end{tabular}
\end{table}

The pipeline stages are configurable: users can set
\texttt{start\_stage} and \texttt{stop\_stage} in YAML to run any
contiguous subset of the pipeline.

% ── 4. Architecture ──────────────────────────────────────────────────────────
\section{Architecture}
\label{sec:architecture}

\subsection{Planner/Designer/Critic Loop}
\label{sec:pdc}

Every stage agent uses the same internal architecture: three sub-agents
coordinated by a planner (Figure~\ref{fig:agents}).

\begin{figure}[H]
  \centering
  \includegraphics[width=\textwidth]{figures/agents_internal.png}
  \caption{Internal agent architecture.  The planner orchestrates the
    designer and critic through tool calls.  Checkpointing enables
    rollback when quality degrades.}
  \label{fig:agents}
\end{figure}

\paragraph{Designer.}
Creates and refines the artifact (plan text or Mermaid diagram).  It
receives the context---user brief or parent plan---produces an initial
design, and iteratively improves it based on critic feedback.  Its full
conversation history is persisted in a SQLite session so the critic and
subsequent iterations have complete context.

\paragraph{Critic.}
Evaluates the current artifact against domain-specific criteria
(structure, clarity, completeness) and returns structured feedback with
per-category scores.  For project, module, and task stages the critic
reads the designer's session history.  For the visualization stage the
critic is a \emph{vision-language model} (VLM) that receives the
rendered diagram image alongside the plan text.

\paragraph{Planner.}
Orchestrates the loop via internal tools:
\begin{itemize}[leftmargin=*]
  \item \texttt{request\_initial\_design} -- ask the designer for the
    first draft.
  \item \texttt{request\_critique} -- ask the critic to evaluate the
    current artifact.
  \item \texttt{request\_design\_change} -- ask the designer to address
    specific issues raised by the critic.
  \item \texttt{reset\_plan\_to\_previous\_checkpoint} -- roll back to
    the last good state if scores degrade.
\end{itemize}

The planner alternates between critique and design-change rounds until
quality thresholds are met or the maximum number of rounds is reached.
\texttt{max\_critique\_rounds} and \texttt{early\_finish\_min\_score}
are configurable per agent.

\subsection{Checkpoint-Based Auto-Reset}
\label{sec:checkpoint}

Each critique round updates a checkpoint consisting of the current
plan state and associated scores.  The framework maintains an
\emph{N-1/N} checkpoint pair: the current checkpoint~(N) and the
previous one~(N-1).  If the total score after a critique is lower than
the previous checkpoint, the planner can call
\texttt{reset\_plan\_to\_previous\_checkpoint}, which:
\begin{enumerate}[leftmargin=*]
  \item Compares the total scores of checkpoints N and N-1.
  \item If N $<$ N-1, restores the plan state from checkpoint N-1.
  \item Appends a rollback message to the designer session so subsequent
    iterations continue from the better state.
\end{enumerate}
This mechanism provides \textbf{intra-run fault tolerance}: the system
recovers from quality regressions without manual intervention.

\subsection{Structured Scoring}
\label{sec:scoring}

Critic output is a structured dataclass (\texttt{CritiqueWithScores})
with domain-specific categories.  Each category has a name, a grade
(0--10), and a textual comment.  Domain-specific subclasses define the
relevant categories:

\begin{itemize}[leftmargin=*]
  \item \texttt{ProjectCritiqueWithScores} -- project-level criteria.
  \item \texttt{ModuleCritiqueWithScores} -- module-level criteria.
  \item \texttt{TaskCritiqueWithScores} -- task-level criteria.
  \item \texttt{VisualizationCritiqueWithScores} -- diagram-level criteria.
\end{itemize}

A shared \texttt{compute\_total\_score} function sums category grades,
and \texttt{format\_score\_deltas\_for\_planner} formats score changes
between iterations so the planner can make informed decisions about
when to continue, stop, or roll back.

\subsection{Inversion of Control and Configuration}
\label{sec:ioc}

\sprintagent{} follows the \ioc{} principle~\cite{fowler2004ioc}: the
framework owns the execution flow and calls user-supplied components at
well-defined extension points.  Configuration is managed through
Hydra~\cite{hydra2019} and OmegaConf~\cite{omegaconf2019}, providing:

\begin{itemize}[leftmargin=*]
  \item \textbf{Hierarchical YAML composition.}  Each agent type has a
    \texttt{core\_*.yaml} file with shared settings and a
    \texttt{workflow\_*.yaml} overlay for workflow-specific overrides.
    Hydra's defaults list composes them automatically.
  \item \textbf{Cross-config interpolation.}  Global settings (e.g.,
    \texttt{openai.service\_tier}) are defined once and interpolated
    into agent configs via \texttt{\$\{openai.service\_tier\}}.
  \item \textbf{Custom resolvers.}  OmegaConf resolvers for path
    resolution (\texttt{resolve\_path}, \texttt{resolve\_list}),
    boolean negation (\texttt{not}), equality checks (\texttt{equal}),
    and conditionals (\texttt{ifelse}) enable expressive config logic
    without code.
  \item \textbf{Agent wiring.}  Each experiment declares a
    \texttt{compatible\_*\_agents} dictionary mapping config names to
    Python classes.  At runtime the framework reads the
    \texttt{\_name} field from the resolved config and instantiates the
    corresponding class, achieving dependency injection through
    configuration.
\end{itemize}

Adding a new agent type requires only a YAML config file and a
registration entry in the experiment's compatible agents dictionary---no
changes to the core orchestration code.

\subsection{Session Management and Turn Trimming}
\label{sec:session}

Designer and critic conversations are persisted in SQLite databases
(e.g., \texttt{prompt\_000\_project\_designer.db}).  This provides:

\begin{itemize}[leftmargin=*]
  \item \textbf{Crash recovery.}  Sessions survive process restarts.
  \item \textbf{Context for critics.}  The critic reads the designer's
    full session history to evaluate the current state in context.
  \item \textbf{Auditability.}  Every interaction is recorded and can be
    inspected post-hoc.
\end{itemize}

For long-running agents, unbounded context growth exceeds token limits.
\sprintagent{} addresses this with \texttt{TurnTrimmingSession}, a
decorator that wraps the base \texttt{SQLiteSession}:

\begin{enumerate}[leftmargin=*]
  \item \textbf{Turn parsing.}  Conversation items are segmented into
    turns by user-message boundaries.
  \item \textbf{Retention window.}  The last $N$ turns (configurable via
    \texttt{keep\_last\_n\_turns}) are kept intact.
  \item \textbf{Image stripping.}  Older turns have images replaced with
    text placeholders to reduce token usage.
  \item \textbf{Optional summarization.}  When
    \texttt{enable\_summarization} is active, trimmed turns are
    summarized by an LLM.  Summaries are cached in a dedicated SQLite
    table (keyed by turn content hash) to avoid redundant API calls.
\end{enumerate}

An additional \texttt{IntraTurnImageFilter} strips images from older
tool outputs \emph{within} a single turn, keeping only the most recent
observations.  This is particularly useful for the visualization critic,
which may call observation tools multiple times per iteration.

\subsection{Parallel Execution and Process Isolation}
\label{sec:parallel}

\sprintagent{} supports both sequential and parallel execution of
prompts.  Parallel execution uses \texttt{ProcessPoolExecutor}, where
each prompt runs in a separate process:

\begin{itemize}[leftmargin=*]
  \item \textbf{Memory isolation.}  Each process has its own address
    space; no shared mutable state.
  \item \textbf{Fault containment.}  A crash or exception in one worker
    does not affect others.  The orchestrator collects results and
    reports per-prompt success or failure.
  \item \textbf{Log separation.}  Each worker creates its own
    \texttt{FileLoggingContext} writing to
    \texttt{prompt\_XXX/plan.log}, with \texttt{suppress\_stdout=True}
    to prevent log interleaving.
\end{itemize}

The number of workers is controlled by the \texttt{num\_workers}
configuration parameter; setting it to~1 enables sequential execution
for debugging.

\subsection{Unified Logging}
\label{sec:logging}

\sprintagent{} uses a \texttt{FileLoggingContext} context manager that
attaches a file handler to Python's root logger, capturing output from
all modules without per-module configuration:

\begin{itemize}[leftmargin=*]
  \item \textbf{Experiment-level log.}  The main process writes to
    \texttt{experiment.log} with both file and console output.
  \item \textbf{Per-prompt logs.}  Parallel workers write to isolated
    \texttt{plan.log} files with console output suppressed.
  \item \textbf{Consistent format.}  All log entries follow
    \texttt{timestamp -- logger\_name -- level -- message}.
\end{itemize}

Combined with the SQLite session databases, this creates a full audit
trail: for every prompt one can inspect the experiment log, the
per-prompt log, and the designer/critic conversation history.

\subsection{Prompt Management}
\label{sec:prompts}

Prompts are stored as YAML files with Jinja2 templates under a
structured directory (\texttt{src/prompts/data/<agent\_type>/}).  A
\texttt{PromptRegistry} provides type-safe access through
domain-specific enums (e.g., \texttt{ProjectAgentPrompts.PLANNER\_AGENT})
and enforces strict template-variable validation: the set of provided
keyword arguments must exactly match the declared
\texttt{template\_variables}, preventing silent errors from missing or
extra parameters.

Prompts are loaded and cached by a \texttt{PromptManager} with
LRU caching, and rendered at runtime with the appropriate context
(e.g., the user brief, a parent plan, or critic feedback).

% ── 5. Use Case ──────────────────────────────────────────────────────────────
\section{Use Case: Agentic Project Management}
\label{sec:usecase}

To demonstrate how the framework looks in practice, an autonomous
project-management pipeline is built on top of \sprintagent{} as a
reference application.  This use case is not the primary contribution;
it serves to illustrate how practitioners can wire their own agents,
tools, and domain logic into the scaffold.  Given a natural-language brief such as
\emph{``Plan the frontend dashboard module, including UI design tasks
and testing tasks,''} the system produces:

\begin{enumerate}[leftmargin=*]
  \item \textbf{Project plan} -- High-level components, boundaries, and
    dependencies (\texttt{project\_plan.md}).
  \item \textbf{Module plan} -- Module structure per component, refined
    from the project plan (\texttt{module\_plan.md}).
  \item \textbf{Task plan} -- Concrete, assignable tasks per module
    (\texttt{task\_plan.md}).
  \item \textbf{Visual diagrams} -- Mermaid-based hierarchy diagrams
    rendered to PNG and PDF for each plan level.
\end{enumerate}

The pipeline is fully autonomous: the user provides a brief, and agents
iteratively refine plans until quality criteria are met.  Each stage
consumes the output of the previous one, ensuring alignment from
high-level strategy down to individual tasks.

The visualization stage is noteworthy because its critic is a
vision-language model.  After the designer produces a Mermaid diagram
and the system renders it to an image, the VLM critic compares the
rendered diagram against the source plan, checking for structural
accuracy, completeness, and readability.  This closes the loop between
textual plans and visual artifacts.

Outputs are organized under
\texttt{outputs/<timestamp>/<run>/prompt\_<id>/}, with all plans,
diagrams, logs, and session databases co-located for easy inspection.

% ── 6. Implementation ────────────────────────────────────────────────────────
\section{Implementation Details}
\label{sec:implementation}

\sprintagent{} is implemented in Python and built on the OpenAI Agents
SDK~\cite{openaiagentssdk2025} for agent execution.  Key technology
choices include:

\begin{table}[H]
  \centering
  \caption{Technology stack.}
  \label{tab:stack}
  \begin{tabular}{@{}ll@{}}
    \toprule
    \textbf{Component}          & \textbf{Technology} \\
    \midrule
    Agent execution             & OpenAI Agents SDK \\
    Configuration               & Hydra + OmegaConf \\
    Session persistence         & SQLite \\
    Diagram rendering           & Mermaid CLI (Node.js) \\
    Parallel execution          & \texttt{ProcessPoolExecutor} \\
    Prompt templating           & Jinja2 + YAML \\
    Package management          & \texttt{uv} \\
    \bottomrule
  \end{tabular}
\end{table}

The codebase is organized as follows:

\begin{itemize}[leftmargin=*]
  \item \texttt{src/agent\_utils/} -- Shared agent infrastructure:
    base stateful agent, checkpointing, scoring, turn trimming, and
    image filtering.
  \item \texttt{src/\{project,module,task,visualization\}\_agents/} --
    Domain-specific agent implementations.
  \item \texttt{src/experiments/} -- Experiment orchestration and
    pipeline execution.
  \item \texttt{src/prompts/} -- Prompt registry, manager, and YAML
    templates.
  \item \texttt{src/utils/} -- Logging, parallel execution, OpenAI
    helpers, and Mermaid rendering.
  \item \texttt{configurations/} -- Hydra YAML configs for agents,
    experiments, and pipelines.
\end{itemize}

% ── 7. Design Patterns ──────────────────────────────────────────────────────
\section{Design Patterns}
\label{sec:patterns}

Table~\ref{tab:patterns} summarizes the software design patterns
employed in \sprintagent{} and their roles.

\begin{table}[H]
  \centering
  \caption{Design patterns in \sprintagent{}.}
  \label{tab:patterns}
  \begin{tabular}{@{}lll@{}}
    \toprule
    \textbf{Pattern} & \textbf{Location} & \textbf{Purpose} \\
    \midrule
    Template Method  & \texttt{BaseStatefulAgent}     & Domain-specific hooks for prompts and behavior \\
    Decorator        & \texttt{TurnTrimmingSession}   & Wraps sessions with trimming/summarization \\
    Strategy         & Prompt registry, scoring       & Pluggable prompts and scoring schemas \\
    Factory Method   & \texttt{\_create\_*\_agent()}   & Agent creation from configuration \\
    \ioc{}           & Hydra + compatible agents dict & Config-driven agent wiring \\
    Filter           & \texttt{IntraTurnImageFilter}  & Pre-call input filtering for token control \\
    Abstract Base    & \texttt{CritiqueWithScores}    & Shared scoring interface across domains \\
    N-1/N Checkpoint & Checkpoint state               & Rollback on quality degradation \\
    \bottomrule
  \end{tabular}
\end{table}

% ── 8. Discussion ────────────────────────────────────────────────────────────
\section{Discussion}
\label{sec:discussion}

As outlined in Section~\ref{sec:introduction}, building agentic systems
demands deliberate engineering trade-offs.  This section revisits each
trade-off and discusses how \sprintagent{} addresses it.

\paragraph{Complexity vs.\ explainability.}
A recurring tension in agentic systems is between the power of complex
multi-agent orchestration and the need for humans to understand what the
system did and why.  \sprintagent{} addresses this by making every
agent interaction auditable (SQLite sessions), every decision logged
(unified logging), and every quality judgment structured (scored
critiques).  The checkpoint mechanism further supports explainability:
when a rollback occurs, the system records what degraded and what was
restored.

\paragraph{Isolation vs.\ throughput.}
Process isolation provides strong fault containment but incurs overhead
from process creation and inter-process serialization (all arguments
must be picklable).  For most agentic workloads, where LLM API latency
dominates, this overhead is negligible.  The framework also supports
sequential execution for debugging and development.

\paragraph{Context growth vs.\ performance.}
Agents accumulate conversation history that, if left unmanaged, exceeds
token limits and inflates API costs.  \sprintagent{} balances context
richness against performance through turn trimming (keeping only the
last $N$ turns intact), image stripping for older observations, and
optional LLM-based summarization with hash-keyed caching to avoid
redundant calls.  The result is bounded context windows without losing
the essential reasoning trace.

\paragraph{Configuration flexibility vs.\ simplicity.}
Real-world agentic systems have dozens of knobs---model selection,
scoring thresholds, memory policies, pipeline stages---and the
configuration system must scale without becoming its own problem.
\sprintagent{} addresses this through Hydra's hierarchical YAML
composition, cross-config interpolation, and custom OmegaConf resolvers.
The \ioc{} approach means that adding a new agent type, experiment, or
scoring schema requires only YAML and a Python class---no changes to the
orchestration core.  This separation of wiring from logic reduces the
cost of experimentation and lowers the barrier for new contributors.

\paragraph{Limitations.}
The current implementation is tightly coupled to the OpenAI Agents SDK
for agent execution.  Supporting alternative LLM providers would require
an adapter layer.  The checkpoint mechanism operates at the plan level
and does not yet support fine-grained rollback of individual tool calls.
Evaluation is currently qualitative; integrating automated benchmarks
for plan quality is left for future work.

% ── 9. Related Work ──────────────────────────────────────────────────────────
\section{Related Work}
\label{sec:related}

\textbf{Self-improving agents.}
ReAct~\cite{yao2023react} and Reflexion~\cite{shinn2023reflexion}
introduced reasoning-and-acting loops and verbal self-reflection.
\sprintagent{}'s planner/designer/critic loop can be seen as a
structured variant of self-reflection, where the critic provides
quantitative scores and the planner makes rollback decisions based on
score trajectories.

\textbf{Configuration and IoC.}
The \ioc{} pattern~\cite{fowler2004ioc} is well-established in
enterprise software but underutilized in LLM-agent frameworks.
Hydra~\cite{hydra2019} and OmegaConf~\cite{omegaconf2019} are widely
used in ML experiment management; \sprintagent{} adapts them for agent
orchestration.

% ── 10. Conclusion ───────────────────────────────────────────────────────────
\section{Conclusion}
\label{sec:conclusion}

This paper presented \sprintagent{}, a production-ready framework for
hierarchical autonomous agentic systems, developed with four core
engineering trade-offs in mind: complexity vs.\ explainability,
isolation vs.\ throughput, context growth vs.\ performance, and
configuration flexibility vs.\ simplicity.  The framework provides
planner/designer/critic orchestration per agent, \ioc{}-based
configuration, SQLite session management with turn trimming, checkpoint
auto-reset on degradation, process-isolated parallel execution, and
unified logging---all driven by YAML configuration.  Its modular design
allows practitioners to plug their own logic into each module and extend
on top.  Through a demonstration application---an agentic
project-management system built on the framework---the paper showed what
it looks like in practice to build a complete hierarchical agentic
application on \sprintagent{}, confirming that the scaffold is both
functional and extensible.

Future work includes supporting alternative LLM backends, adding
quantitative benchmarks for plan quality, extending the checkpoint
mechanism to sub-agent granularity, and exploring distributed execution
beyond single-machine process pools.

\sprintagent{} is open source and available at
\url{https://github.com/alishams21/sprintagent}.

% ── Bibliography ─────────────────────────────────────────────────────────────
\bibliographystyle{plainnat}
\bibliography{references}

\end{document}
